%============================================
% Fixed Pearson-style STEM Book Template
% with working tcolorbox/marginnote definitions
%============================================
\documentclass[12pt,oneside]{book}
\usepackage[utf8]{inputenc}
\usepackage[T1]{fontenc}
\usepackage{lmodern}        % change if desired e.g. lmodern

%------------------ User-configurable macros ------------------
\newcommand{\BookTitle}{My STEM Book Title}
\newcommand{\BookSubtitle}{An Informative Subtitle}
\newcommand{\AuthorName}{Author Name}
\newcommand{\RevisionDate}{September~2025}



%------------------ Packages & Layout ------------------


% macros
\newcommand{\TopMg}{1.0in}
\newcommand{\InnerMg}{1.0in}\newcommand{\OutMg}{1.0in}
\newcommand{\BottomMg}{1.0in}
\usepackage[
  twoside, 
  % includemp,                        % <--- include margin notes in layout calculation
%   showframe,                        % <--- uncomment for design
  letterpaper
]{geometry}
\geometry{
  top=\TopMg,
  inner=\InnerMg,outer=\OutMg,
  bottom=\BottomMg,
}




\usepackage{microtype}
\usepackage{amsmath,amssymb,amsthm}
\usepackage{mathtools}
\usepackage{graphicx}
\usepackage{subcaption}
\usepackage{booktabs}
\usepackage{caption}
\usepackage{xcolor}
\usepackage{enumitem}


\usepackage{fancyhdr}
\pagestyle{fancy}
\fancyhf{}\renewcommand{\headrulewidth}{0pt}\fancyfoot[C]{\thepage}










% tcolorbox
% Important: load tcolorbox once with libraries (most) so 'breakable' etc. work.
% See tcolorbox manual / CTAN for recommended usage. :contentReference[oaicite:1]{index=1}
\usepackage[most]{tcolorbox}

% ---------------- Box definitions (after tcolorbox is loaded) ----------------
% Use a tcolorbox style (pgfkeys) instead of a macro to avoid '#'
% parameter issues inside key definitions.
\tcbset{tcbbase/.style={
  enhanced, breakable,
  left=4pt, right=4pt, top=3pt, bottom=3pt,
  fonttitle=\bfseries
}}

% Sidebar: customizable title via optional arg
\newtcolorbox{sidebarbox}[1][]{tcbbase,
  colback=gray!10, colframe=gray!60, title={#1}}

% Worked example: slightly wider padding and prefixed title
\newtcolorbox{workedexample}[1][]{tcbbase,
  left=6pt, right=6pt, top=6pt, bottom=6pt,
  colback=blue!5, colframe=blue!60,
  title={Worked Example: #1}}

% Tip/Info/Warning fixed titles with colors
\newtcolorbox{TipBox}{tcbbase,
  colback=green!5, colframe=green!50, title=Tip}
\newtcolorbox{InfoBox}{tcbbase,
  colback=cyan!5, colframe=cyan!60, title=Info}
\newtcolorbox{WarningBox}{tcbbase,
  colback=red!10, colframe=red!60, title=Warning}

% Theorem / definition environments
\newtheorem{theorem}{Theorem}[chapter]
\newtheorem{lemma}[theorem]{Lemma}
\theoremstyle{definition}
\newtheorem{definition}[theorem]{Definition}
\newtheorem{example}[theorem]{Example}
\theoremstyle{remark}
\newtheorem{remark}[theorem]{Remark}
\newtheorem{exercise}[theorem]{Exercise}


%------------------ Document content ------------------
\begin{document}
\frontmatter
\title{\BookTitle\\[1ex]\large\BookSubtitle}
\author{\AuthorName}
\date{\RevisionDate}
\maketitle

\chapter*{Preface}
% Preface text here

\tableofcontents
\listoffigures
\listoftables







\mainmatter
\chapter{Introduction}
\section{Background}

This is sample text.


\begin{sidebarbox}[Historical Note]
Sidebars should be short. They are implemented using tcolorbox and are breakable.
\end{sidebarbox}

\section{Definitions}
\begin{definition}
A *ring* \(R\) is a set with two binary operations ...
\end{definition}

\begin{workedexample}[Ideal in $\mathbb{Z}$]
Let \(I=2\mathbb{Z}\). Show $I$ is an ideal: closure and absorption by ring elements.
\end{workedexample}

\begin{TipBox}
Always define notation clearly.
\end{TipBox}
\begin{InfoBox}
The ring $\mathbb{Z}/n\mathbb{Z}$ has exactly $n$ elements.
\end{InfoBox}
\begin{WarningBox}
A ring need not be a field (invertibility is not automatic).
\end{WarningBox}

\section{Exercises}
\begin{exercise}
Prove that the intersection of two ideals is an ideal.
\end{exercise}



\begin{sidebarbox}[Historical Note]
Sidebars should be short. They are implemented using tcolorbox and are breakable.
\end{sidebarbox}




\section{Definitions}
\begin{definition}
A *ring* \(R\) is a set with two binary operations ...
\end{definition}

\begin{workedexample}[Ideal in $\mathbb{Z}$]
Let \(I=2\mathbb{Z}\). Show $I$ is an ideal: closure and absorption by ring elements.
\end{workedexample}



\end{document}